%% kapa: exam=1 round=36 year=2025 period=2 form=A subject=appraisal_law subject_ko=감정평가관계법규 seq=1 number=1
%% source: Q-Net 감정평가사 기출문제 (cst00309) artlSeq=5231525 file=raw/1008/5231525/2025년 제36회 감정평가사 1차시험 2교시 문제지 원본.pdf
%% conversion: kordoc 4.24.51 -> tools/gen1.py (auto); review=proofread-sonnet; flags=none
\begin{question}{1}
국토의 계획 및 이용에 관한 법령상 광역도시계획에 관한 설명으로 옳지 않은 것은?
\begin{choices}
  \item 광역도시계획은 광역계획권의 장기발전방향을 제시하는 계획을 말한다.
  \item 국토교통부장관은 광역계획권을 지정하려면 관계 시ㆍ도지사, 시장 또는 군수의 의견을 들은 후 중앙도시계획위원회의 심의를 거쳐야 한다.
  \item 국토교통부장관은 시ㆍ도지사가 요청하는 경우에는 관할 시ㆍ도지사와 공동으로 광역도시계획을 수립할 수 있다.
  \item 도지사는 시장 또는 군수가 협의를 거쳐 요청하는 경우에는 단독으로 광역도시계획을 수립할 수 있다.
  \item 시장 또는 군수는 광역도시계획을 수립하려면 국토교통부장관의 승인을 받아야 한다.
\end{choices}
\end{question}
\answer{5}
